\chapter[13]{Qualitative Behavior of Classes of Solutions}
\label{ch:chow-ii-qualitative-behavior}

A parallel-transport invariant family \(\Gamma(x)\) of subsets of the
algebraic curvature operators is \emph{preserved} if
\(\Rm(g(0))\in\Gamma\) implies \(\Rm(g(t))\in\Gamma\) throughout the flow.
It is \emph{attractive} if every complete pointed singularity model formed at
a finite-time singularity has curvature in \(\Gamma\) at its base time.

\section{Curvature conditions that are not preserved}

\subsection{Examples of preserved curvature cones}

\begin{example}[Nonnegative scalar curvature]
\label{ex:chow-ii-nonnegative-scalar-preserved-attractive}
\source{chow-et-al-ricci-flow-part-ii:page-0225}
\dcref{ch13:13.1}
\group{Preserved Curvature Conditions}

For every Ricci flow to which the scalar maximum principle applies,
\(R(\cdot,t)\geq\min R(\cdot,0)\).  Nonnegative scalar curvature is both
preserved and attractive.
\end{example}
\begin{proof}
\sketch
\uses{thm:chow-ii-scalar-wmp}
Since \(\partial_tR=\Delta R+2|\Ric|^2\geq\Delta R\), the maximum principle
gives the lower bound.  Under a blow-up by \(Q_i\to\infty\), it becomes
\(R_i\geq Q_i^{-1}\min R(0)\); hence every smooth limit has \(R\geq0\).
\end{proof}

\begin{example}[Nonnegative curvature operator]
\label{ex:chow-ii-nonnegative-curvature-operator-preserved}
\source{chow-et-al-ricci-flow-part-ii:page-0225}
\dcref{ch13:13.2}
\group{Preserved Curvature Conditions}

Nonnegative curvature operator is preserved by Ricci flow in every dimension.
In dimensions two and three it is equivalent to nonnegative sectional
curvature.  Positive isotropic curvature is also preserved in every dimension.
\end{example}
\begin{proof}
\sketch
\uses{thm:chow-ii-curvature-operator-preservation}
The relevant cones are closed, convex, invariant under parallel transport,
and invariant under the curvature reaction ODE.  The tensor maximum principle
preserves them.  In dimensions at most three every two-form is decomposable,
which gives the stated equivalence.
\end{proof}

\begin{example}[Hamilton--Ivey estimate]
\label{ex:chow-ii-hamilton-ivey-attractiveness}
\source{chow-et-al-ricci-flow-part-ii:page-0225}
\dcref{ch13:13.3}
\group{Preserved Curvature Conditions}

For a three-dimensional Ricci flow normalized as in the Hamilton--Ivey
estimate, if \(\lambda_1<0\) is the least curvature-operator eigenvalue, then
\[
 R(x,t)\geq|\lambda_1|(\log|\lambda_1|+\log(1+t)-3).
\]
Consequently nonnegative curvature operator is attractive.
\end{example}
\begin{proof}
\sketch
\uses{thm:chow-ii-hamilton-ivey}
The inequality is the Hamilton--Ivey estimate.  A negative eigenvalue in a
smooth blow-up limit would make the logarithmic right-hand side, after
rescaling, unbounded, contradicting smooth convergence.
\end{proof}

\begin{example}[Two-nonnegative curvature operator]
\label{ex:chow-ii-two-nonnegative-curvature-preserved}
\source{chow-et-al-ricci-flow-part-ii:page-0225,chow-et-al-ricci-flow-part-ii:page-0226}
\dcref{ch13:13.4}
\group{Preserved Curvature Conditions}

If \(\lambda_1\leq\cdots\leq\lambda_N\) are the eigenvalues of an algebraic
curvature operator, then \(\lambda_1+\lambda_2\geq0\) is preserved by Ricci
flow.  In dimension three it is equivalent to nonnegative Ricci curvature;
strict positivity forces a spherical space form.
\end{example}
\begin{proof}
\sketch
\uses{prop:chow-ii-ricci-flow-preserves-two-nonnegative,thm:chow-ii-boehm-wilking-convergence}
The first assertion is the two-nonnegative curvature maximum principle.  In
dimension three the Ricci eigenvalues are the pairwise curvature-operator
eigenvalue sums.  Under strict positivity, normalized flow converges to
constant positive curvature, so the manifold is a spherical space form.
\end{proof}

\begin{example}[Nonnegative holomorphic bisectional curvature]
\label{ex:chow-ii-bisectional-curvature-preserved}
\source{chow-et-al-ricci-flow-part-ii:page-0226}
\dcref{ch13:13.5}
\group{Kahler Curvature Conditions}

Nonnegative holomorphic bisectional curvature is preserved by Kähler--Ricci
flow whenever the maximum principle applies.
\end{example}
\begin{proof}
\sketch
\uses{thm:chow-ii-system-wmp}
Its Kähler curvature cone is closed, convex, unitary-invariant, and preserved
by the reaction ODE, so the bundle maximum principle applies.
\end{proof}

\begin{example}[Two-nonnegative traceless Kähler curvature]
\label{ex:chow-ii-traceless-kahler-curvature-preserved}
\source{chow-et-al-ricci-flow-part-ii:page-0226}
\dcref{ch13:13.6}
\group{Kahler Curvature Conditions}

If initially the Ricci tensor is nonnegative and the sum of any two
eigenvalues of the traceless curvature operator on traceless \((1,1)\)-forms
is nonnegative, both conditions are preserved by Kähler--Ricci flow.
\end{example}
\begin{proof}
\sketch
\uses{thm:chow-ii-system-wmp}
The two inequalities form a closed convex unitary-invariant cone, and the
Kähler curvature reaction ODE points into its tangent cone at each boundary
point.  The bundle maximum principle gives preservation.
\end{proof}

\begin{example}[Orthogonal bisectional curvature preserves Ricci curvature]
\label{ex:chow-ii-orthogonal-bisectional-preserves-ricci}
\source{chow-et-al-ricci-flow-part-ii:page-0226}
\dcref{ch13:13.7}
\group{Kahler Curvature Conditions}

If a Kähler--Ricci flow starts with nonnegative Ricci curvature and positive
orthogonal bisectional curvature is preserved, then Ricci curvature remains
nonnegative.
\end{example}
\begin{proof}
\sketch
At a first zero of the least Ricci eigenvalue, its reaction term is a sum of
orthogonal bisectional curvatures weighted by nonnegative Ricci eigenvalues.
It and the Laplacian term are nonnegative, so the maximum principle excludes
passage into negative Ricci curvature.
\end{proof}

\subsection{Riemannian submersions and sectional curvature}

For a Riemannian submersion \(\pi:(\mathcal X,G)\to(\mathcal Y,g)\), with
vertical and horizontal projections \(\mathcal V,\mathcal H\), set
\[
 A_\varphi\psi=\mathcal H\nabla_{\mathcal H\varphi}\mathcal V\psi
 +\mathcal V\nabla_{\mathcal H\varphi}\mathcal H\psi,
 \quad
 T_\varphi\psi=\mathcal H\nabla_{\mathcal V\varphi}\mathcal V\psi
 +\mathcal V\nabla_{\mathcal V\varphi}\mathcal H\psi.
\]

\begin{example}[Symmetries of the submersion tensors]
\label{ex:chow-ii-oneill-tensor-symmetries}
\source{chow-et-al-ricci-flow-part-ii:page-0227}
\dcref{ch13:13.8}
\group{Riemannian Submersions}

For vertical \(U,V\) and horizontal \(X,Y\),
\[
 A_XY=-A_YX,\quad \langle A_XU,Y\rangle=-\langle A_XY,U\rangle,
\quad T_UV=T_VU,\quad
\langle T_UX,V\rangle=-\langle T_UV,X\rangle.
\]
\end{example}
\begin{proof}
\sketch
Metric compatibility gives the inner-product identities.  Integrability of
the vertical distribution and torsion-freeness give \(T_UV=T_VU\).  The
Koszul formula gives \(A_XY=\tfrac12\mathcal V[X,Y]\), hence skew-symmetry.
\end{proof}

\begin{theorem}[Curvatures of a Riemannian submersion]
\label{thm:chow-ii-oneill-curvature-formulas}
\source{chow-et-al-ricci-flow-part-ii:page-0227}
\dcref{ch13:13.9}
\group{Riemannian Submersions}

For vertical \(U,V\) and horizontal \(X,Y\),
\[
 K_{\mathcal X}(U,V)=K_{\mathcal Z}(U,V)
 +\frac{|T_UV|^2-\langle T_UU,T_VV\rangle}{|U\wedge V|^2},
\]
\[
 K_{\mathcal X}(U,X)=
 \frac{|A_XU|^2+\langle(\nabla_XT)_UU,X\rangle-|T_UX|^2}{|U\wedge X|^2},
\]
\[
 K_{\mathcal X}(X,Y)=K_{\mathcal Y}(\pi_*X,\pi_*Y)
 -3\frac{|A_XY|^2}{|X\wedge Y|^2}.
\]
\end{theorem}
\begin{proof}
\sketch
\uses{ex:chow-ii-oneill-tensor-symmetries}
Decompose all covariant derivatives in the curvature definition into vertical
and horizontal parts.  The four tensor identities reduce the vertical case
to the fiber Gauss equation, the mixed case to the Codazzi equation, and the
horizontal case to
\(\langle R^{\mathcal X}(X,Y)Y,X\rangle
=\langle R^{\mathcal Y}(\pi_*X,\pi_*Y)\pi_*Y,\pi_*X\rangle-3|A_XY|^2\).
Dividing by plane areas proves the formulas.
\end{proof}

\begin{remark}[A submersion increases horizontal curvature]
\label{rem:chow-ii-submersion-increases-curvature}
\source{chow-et-al-ricci-flow-part-ii:page-0227}
\uses{thm:chow-ii-oneill-curvature-formulas}
\dcref{ch13:13.10}
\group{Riemannian Submersions}


The curvature of a plane in the base is at least that of any horizontal plane
mapping isometrically onto it.
\end{remark}

\begin{remark}[Submersions with totally geodesic fibers]
\label{rem:chow-ii-oneill-totally-geodesic-fibers}
\source{chow-et-al-ricci-flow-part-ii:page-0228}
\uses{thm:chow-ii-oneill-curvature-formulas}
\dcref{ch13:13.11}
\group{Riemannian Submersions}


If every fiber is totally geodesic, then \(T=0\) and
\[
 K_{\mathcal X}(U,V)=K_{\mathcal Z}(U,V),\quad
 K_{\mathcal X}(U,X)=\frac{|A_XU|^2}{|U\wedge X|^2},
\]
\[
 K_{\mathcal X}(X,Y)=K_{\mathcal Y}(\pi_*X,\pi_*Y)
 -3\frac{|A_XY|^2}{|X\wedge Y|^2}.
\]
\end{remark}

\subsection{Nonnegative sectional curvature is not preserved}

Let
\(\mathcal Y^4=(\operatorname{SO}(3)\times\mathbb R^2)/\operatorname{SO}(2)
\cong TS^2\) carry the quotient metric.  The submersion to \(S^2\) has
totally geodesic fibers with metric
\(dr^2+r^2(1+r^2)^{-1}d\theta^2\); the preceding formulas show that the
complete metric has bounded nonnegative sectional curvature.

\begin{theorem}[Splitting under persistent nonnegative sectional curvature]
\label{thm:chow-ii-persistent-nonnegative-sectional-splitting}
\source{chow-et-al-ricci-flow-part-ii:page-0230,chow-et-al-ricci-flow-part-ii:page-0231}
\dcref{ch13:13.12}
\group{Unpreserved Curvature Conditions}

Let \((\mathcal M^n,g(t))\), \(0\leq t\leq\delta\), have uniformly bounded
nonnegative sectional curvature, and let \(\widetilde{\mathcal M}\) be its
universal cover.  Then
\[
 (\widetilde{\mathcal M},\widetilde g(t))
 =(\mathcal N_1^{n-k},g_1(t))\times(\mathcal N_2^k,g_2(t)),
\]
where \(\mathcal N_1\) is compact and nonnegatively curved and
\(\mathcal N_2\cong\mathbb R^k\).  For \(t>0\), \(\mathcal N_2\) has a
strictly convex exhaustion and a one-point soul \(p\); the soul of the cover
is \(\mathcal N_1\times\{p\}\).
\end{theorem}
\begin{proof}
\sketch
\uses{thm:system-null-space-rigidity}
Evolve a proper Lipschitz convex Busemann exhaustion on the cover by the heat
equation.  Convexity and properness persist.  At positive time the null space
of its Hessian is parallel by null-space rigidity, so De Rham splitting gives
the stated factors.  Properness makes the null factor compact; strict
convexity on the other factor gives a unique minimum, a point soul, and a
diffeomorphism with \(\mathbb R^k\).  Smooth dependence extends the parallel
splitting to \(t=0\), and the product formula for souls finishes the proof.
\end{proof}

\begin{corollary}[Immediate loss of nonnegative sectional curvature]
\label{cor:chow-ii-sectionally-nonnegative-quotient-counterexample}
\source{chow-et-al-ricci-flow-part-ii:page-0231}
\dcref{ch13:13.13}
\group{Unpreserved Curvature Conditions}

The bounded-curvature Ricci flow starting from the quotient metric on
\(\mathcal Y^4\cong TS^2\) has sectional curvatures of mixed sign at every
positive time.
\end{corollary}
\begin{proof}
\sketch
\uses{thm:chow-ii-persistent-nonnegative-sectional-splitting,rem:chow-ii-oneill-totally-geodesic-fibers}
Otherwise the simply connected \(TS^2\) would be diffeomorphic to
\(N^{4-k}\times\mathbb R^k\), which would trivialize the tangent bundle of
the even-dimensional sphere, contradicting its nonzero Euler class.
\end{proof}

\begin{remark}[Higher-dimensional tangent-bundle counterexamples]
\label{rem:chow-ii-higher-tangent-bundle-counterexamples}
\source{chow-et-al-ricci-flow-part-ii:page-0231}
\uses{thm:chow-ii-persistent-nonnegative-sectional-splitting}
\dcref{ch13:13.14}
\group{Unpreserved Curvature Conditions}


The construction extends to
\((\operatorname{SO}(m+1)\times\mathbb R^m)/\operatorname{SO}(m)\cong TS^m\)
for every \(m\geq2\).  If \(TS^m\) is topologically trivial, the quotient
metric is nevertheless not locally a Riemannian product, again contradicting
the splitting theorem.
\end{remark}

\subsection{Nonnegative Ricci curvature under Kähler--Ricci flow}

\begin{theorem}[Immediate loss of nonnegative Ricci curvature]
\label{thm:chow-ii-kahler-ricci-nonnegative-counterexample}
\source{chow-et-al-ricci-flow-part-ii:page-0232,chow-et-al-ricci-flow-part-ii:page-0233}
\dcref{ch13:13.15}
\group{Unpreserved Curvature Conditions}

For every \(n\geq2\) and \(1\leq k\leq n-1\), the line bundle
\(\mathbb C\hookrightarrow L_k^n\to\mathbb{CP}^{n-1}\) admits a family of
complete bounded-curvature Kähler metrics with nonnegative Ricci curvature
whose Kähler--Ricci flows immediately acquire mixed Ricci curvature.
\end{theorem}
\begin{proof}
\sketch
Write \(r=\log\sum|z_i|^2\), let \(P=P(r)\), and set
\(\varphi=P_r=\sqrt{2e^{kr}/k+m}\), \(m>0\).  Then
\(\varphi\varphi_r=e^{kr}\), and
\(\psi=n-\varphi_r/\varphi-\varphi_{rr}/\varphi_r=n-k\).  The two Ricci
eigenvalues are \(0\) and \((n-k)/\varphi>0\).  Calabi's criterion extends
the metric across the zero section; its asymptotics give completeness and
bounded curvature.  Along the flow,
\[
 \left.\partial_t\psi_r\right|_{t=0}
 =(n-k)\frac{e^{kr}}{\varphi^5}(e^{kr}-km)<0
\]
when \(r<k^{-1}\log(km)\).  Thus the radial eigenvalue becomes negative
immediately while a tangential eigenvalue remains positive.
\end{proof}

\subsection{Compact homogeneous counterexamples}

\begin{theorem}[Böhm--Wilking compact counterexample]
\label{thm:chow-ii-bohm-wilking-ricci-sign-counterexample}
\source{chow-et-al-ricci-flow-part-ii:page-0234,chow-et-al-ricci-flow-part-ii:page-0235,chow-et-al-ricci-flow-part-ii:page-0236,chow-et-al-ricci-flow-part-ii:page-0237}
\dcref{ch13:13.17}
\group{Unpreserved Curvature Conditions}

There are homogeneous metrics with positive sectional curvature on
\(\operatorname{Sp}(3)/(\operatorname{Sp}(1)^3)\) whose Ricci flows
eventually have Ricci curvature of mixed sign.
\end{theorem}
\begin{proof}
\sketch
Write the invariant metric as \(g=x_1Q|_{\mathfrak m_1}\oplus
x_2Q|_{\mathfrak m_2}\oplus x_3Q|_{\mathfrak m_3}\), normalize
\(x_1x_2x_3=1\), and put \(\varphi=x_1+x_2\), \(\psi=x_1-x_2\).  For large
\(\varphi\), small negative \(\psi\), and \(1<c<2\), the normalized flow has
\[
 \varphi'\geq1,\qquad\psi'>0,\qquad
 \frac{d|\psi|}{d\varphi}\geq-c\frac{|\psi|}{\varphi}.
\]
Such initial data may have positive sectional curvature.  The first Ricci
block is
\[
 x_1r_1=16-16\frac{\varphi+\psi}{(\varphi^2-\psi^2)^2}
 +(\varphi+\psi)\varphi\psi.
\]
Integration gives
\((\varphi+\psi)\varphi\psi\leq-\varepsilon\varphi^{2-c}\), so this block
tends to \(-\infty\), whereas scalar curvature stays positive.  Thus Ricci
curvature becomes mixed.  Space-time rescaling transfers the result from
normalized to unnormalized Ricci flow.
\end{proof}

\section{Real analyticity in the space variables}

\subsection{Flat heat equations}

\begin{example}[Criterion for real analyticity]
\label{ex:chow-ii-flat-factorial-analyticity-criterion}
\source{chow-et-al-ricci-flow-part-ii:page-0239,chow-et-al-ricci-flow-part-ii:page-0240,chow-et-al-ricci-flow-part-ii:page-0241}
\dcref{ch13:13.18}
\group{Real Analyticity}

Let \((\mathcal N^n,h)\) be complete and flat.  A smooth function
\(f:\mathcal N\to\mathbb R\) is real analytic if, near every \(a\), there
are \(r,C>0\) such that
\[
 \sum_{m=0}^{\infty}\frac{r^m}{m!}|\nabla^mf|^2(x)\leq C
 \qquad(x\in B(a,r)).
\]
\end{example}
\begin{proof}
\sketch
In a Euclidean chart, Taylor's formula on the segment from \(a\) to \(x\)
expresses the remainder through derivatives of total order \(m+1\).
Cauchy--Schwarz, the series bound, and
\((k_1+\cdots+k_n)!/(k_1!\cdots k_n!)\leq n^{k_1+\cdots+k_n}\) bound it on
\(B(a,r/2)\) by \(C_1 2^{-(m+1)}\).  The remainders converge uniformly to
zero, so the Taylor series converges to \(f\).
\end{proof}

\begin{proposition}[Real analyticity of solutions to the heat equation]
\label{prop:chow-ii-flat-heat-real-analyticity}
\source{chow-et-al-ricci-flow-part-ii:page-0241,chow-et-al-ricci-flow-part-ii:page-0242}
\dcref{ch13:13.19}
\group{Real Analyticity}

Let \((\mathcal N^n,h)\) be closed and flat, and let
\(\partial_tf=\Delta_hf\).  Then
\[
 \sum_{k=0}^{\infty}\frac{(2t)^k}{k!}|\nabla^kf|^2(x,t)
 \leq\max_{\mathcal N}|f(\cdot,0)|^2.
\]
Thus \(f(\cdot,t)\) is real analytic for every \(t>0\).
\end{proposition}
\begin{proof}
\sketch
\uses{ex:chow-ii-flat-factorial-analyticity-criterion}
Flatness gives
\((\partial_t-\Delta)|\nabla^kf|^2=-2|\nabla^{k+1}f|^2\).  Hence
\[
 (\partial_t-\Delta)\sum_{k=0}^{\ell}
 \frac{(2t)^k}{k!}|\nabla^kf|^2
 =-2\frac{(2t)^\ell}{\ell!}|\nabla^{\ell+1}f|^2\leq0.
\]
The maximum principle, followed by \(\ell\to\infty\), proves the estimate;
the preceding criterion proves analyticity.
\end{proof}

\subsection{Ricci-flow analyticity}

\begin{lemma}[Normal-coordinate criterion for an analytic metric]
\label{lem:chow-ii-normal-coordinate-analytic-atlas}
\source{chow-et-al-ricci-flow-part-ii:page-0242}
\dcref{ch13:13.20}
\group{Real Analyticity}

If a Riemannian manifold has a covering subatlas of normal coordinates in
which all metric components are real analytic, then those charts define a
real analytic structure and the metric is real analytic in that structure.
\end{lemma}
\begin{proof}
\sketch
The Christoffel symbols are analytic functions of the metric and its first
derivatives.  The analytic geodesic ODE therefore has solutions depending
analytically on their initial data.  Transition maps between normal charts,
being compositions of exponential maps and their local inverses, are
analytic.  Thus the charts form the required atlas.
\end{proof}

\begin{theorem}[Real analyticity in space of Ricci-flow solutions]
\label{thm:chow-ii-ricci-flow-spatial-real-analyticity}
\source{chow-et-al-ricci-flow-part-ii:page-0242,chow-et-al-ricci-flow-part-ii:page-0243,chow-et-al-ricci-flow-part-ii:page-0244,chow-et-al-ricci-flow-part-ii:page-0245,chow-et-al-ricci-flow-part-ii:page-0246,chow-et-al-ricci-flow-part-ii:page-0247,chow-et-al-ricci-flow-part-ii:page-0248,chow-et-al-ricci-flow-part-ii:page-0249,chow-et-al-ricci-flow-part-ii:page-0250,chow-et-al-ricci-flow-part-ii:page-0251,chow-et-al-ricci-flow-part-ii:page-0252}
\dcref{ch13:13.21}
\group{Real Analyticity}

Let \((\mathcal M^n,g(t))\), \(0\leq t<T\), be a complete Ricci flow with
bounded curvature.  For every fixed \(t\in(0,T)\), \(g(t)\) is real analytic
in the spatial variables.
\end{theorem}
\begin{proof}
\sketch
\uses{lem:chow-ii-normal-coordinate-analytic-atlas,lem:chow-ii-main-factorial-curvature-estimate,lem:chow-ii-normal-coordinate-metric-derivatives,rem:chow-ii-noncompact-analyticity}
Fix \(t_0>0\) and restart at a slightly earlier positive time.  The factorial
curvature estimate gives \(|\nabla^k\Rm(g(t_0))|\leq Ck!r^{-k-2}\).  The
normal-coordinate derivative lemma gives factorial bounds for all coordinate
derivatives of \(g_{ij}(t_0)\), so their Taylor series converge.  The
normal-coordinate criterion supplies the analytic atlas.  The noncompact
case is justified by Remark~\ref{rem:chow-ii-noncompact-analyticity}.
\end{proof}

For tensors \(A,B\), write \(A*B\) for a contraction of \(A\otimes B\).
The relation \(S\preccurlyeq C A*B\) means that \(S\) is a sum of at most
\(C\) such contractions.

\begin{lemma}[Commutator of one covariant derivative and the Laplacian]
\label{lem:chow-ii-nabla-laplacian-commutator}
\source{chow-et-al-ricci-flow-part-ii:page-0243}
\dcref{ch13:13.22}
\group{Curvature Derivative Estimates}

For every covariant \(\ell\)-tensor \(A\),
\[
 (\nabla\Delta-\Delta\nabla)A
 \preccurlyeq n(2\ell+1)\Rm*\nabla A+n\ell\nabla\Rm*A.
\]
\end{lemma}
\begin{proof}
\sketch
In a normal orthonormal frame, commute \(\nabla_i\) past the two traced
derivatives.  The first interchange gives \(\ell+1\) curvature actions on
\(\nabla A\); the second gives another \(\ell\) such actions and \(\ell\)
derivatives of curvature acting on \(A\).  Summing the traced index gives the
stated count.
\end{proof}

\begin{remark}[Second-Bianchi form of the commutator]
\label{rem:chow-ii-nabla-laplacian-bianchi-form}
\source{chow-et-al-ricci-flow-part-ii:page-0243}
\uses{lem:chow-ii-nabla-laplacian-commutator}
\dcref{ch13:13.23}
\group{Curvature Derivative Estimates}


The second Bianchi identity rewrites the derivative-curvature term as
\[
 (\nabla\Delta-\Delta\nabla)A
 \preccurlyeq n(2\ell+1)\Rm*\nabla A+2\ell\nabla\Ric*A.
\]
\end{remark}

\begin{lemma}[Commutator of higher covariant derivatives and the Laplacian]
\label{lem:chow-ii-higher-nabla-laplacian-commutator}
\source{chow-et-al-ricci-flow-part-ii:page-0243,chow-et-al-ricci-flow-part-ii:page-0244,chow-et-al-ricci-flow-part-ii:page-0245}
\dcref{ch13:13.24}
\group{Curvature Derivative Estimates}

If \(T\) is a covariant \(p\)-tensor, \(p\geq1\), then
\[
 \nabla^k\Delta T-\Delta\nabla^kT
 \preccurlyeq n(p+1)\frac{(k+2)!}{2!k!}\Rm*\nabla^kT
\]
\[
 {}+n(2p+1)\sum_{i=1}^k
 \frac{(k+2)!}{(i+2)!(k-i)!}\nabla^i\Rm*\nabla^{k-i}T.
\]
\end{lemma}
\begin{proof}
\sketch
\uses{lem:chow-ii-nabla-laplacian-commutator,rem:chow-ii-differentiable-contraction-count}
Induct on \(k\).  Differentiate the inductive estimate and add the
one-derivative commutator on \(\nabla^kT\).  Leibniz' rule splits each old
summand into adjacent indices; Pascal's identity combines their factorial
coefficients.  The two boundary coefficients are bounded by the corresponding
coefficients in the formula with \(k\) replaced by \(k+1\).
\end{proof}

\begin{remark}[Differentiating contraction-count estimates]
\label{rem:chow-ii-differentiable-contraction-count}
\source{chow-et-al-ricci-flow-part-ii:page-0244}
\dcref{ch13:13.25}
\group{Curvature Derivative Estimates}

If \(S\preccurlyeq C A*B\), then
\(\nabla S\preccurlyeq C\nabla A*B+C A*\nabla B\).  Thus, unlike a norm
inequality, a contraction-count relation may be differentiated.
\end{remark}

\begin{lemma}[Commutator of covariant and time derivatives]
\label{lem:chow-ii-nabla-time-commutator}
\source{chow-et-al-ricci-flow-part-ii:page-0245,chow-et-al-ricci-flow-part-ii:page-0246}
\dcref{ch13:13.26}
\group{Curvature Derivative Estimates}

If \(T\) is a covariant \(p\)-tensor along Ricci flow, then
\[
 \nabla^k\partial_tT-\partial_t\nabla^kT
 \preccurlyeq3n(p+1)\sum_{i=1}^k
 \frac{(k+1)!}{(i+1)!(k-i)!}\nabla^i\Rm*\nabla^{k-i}T.
\]
\end{lemma}
\begin{proof}
\sketch
For \(k=1\), variation of the connection gives
\((\nabla\partial_t-\partial_t\nabla)T=3p\nabla\Ric*T\).  Differentiate
inductively.  Leibniz' rule and Pascal's identity combine adjacent
coefficients, while \(k+p\leq(p+1)(k+1)\) controls the new boundary term.
\end{proof}

\begin{corollary}[Commutator with the heat operator]
\label{cor:chow-ii-nabla-heat-commutator}
\source{chow-et-al-ricci-flow-part-ii:page-0247}
\dcref{ch13:13.27}
\group{Curvature Derivative Estimates}

If \(T\) is a covariant \(p\)-tensor, then
\[
 \nabla^k(\partial_t-\Delta)T-(\partial_t-\Delta)\nabla^kT
 \preccurlyeq3n(p+1)\sum_{i=0}^k(i+k+4)
 \frac{(k+1)!}{(i+2)!(k-i)!}\nabla^i\Rm*\nabla^{k-i}T.
\]
\end{corollary}
\begin{proof}
\sketch
\uses{lem:chow-ii-higher-nabla-laplacian-commutator,lem:chow-ii-nabla-time-commutator}
Subtract the two commutator formulas and put their coefficients over the
common denominator \((i+2)!(k-i)!\).  Their sum is bounded by the displayed
coefficient; the \(i=0\) term is the leading Laplacian term.
\end{proof}

\begin{lemma}[Evolution inequality for curvature derivatives]
\label{lem:chow-ii-curvature-derivative-evolution-inequality}
\source{chow-et-al-ricci-flow-part-ii:page-0248}
\dcref{ch13:13.28}
\group{Curvature Derivative Estimates}

There is \(C_n<\infty\) such that
\[
 \partial_t|\nabla^k\Rm|^2\leq\Delta|\nabla^k\Rm|^2
 -2|\nabla^{k+1}\Rm|^2
\]
\[
 {}+C_n\sum_{i=0}^k\frac{(k+2)!}{(i+2)!(k-i)!}
 |\nabla^i\Rm|\,|\nabla^{k-i}\Rm|\,|\nabla^k\Rm|.
\]
\end{lemma}
\begin{proof}
\sketch
\uses{cor:chow-ii-nabla-heat-commutator}
Apply the commutator estimate to
\((\partial_t-\Delta)\Rm=(4n+5)\Rm*\Rm\), expand by Leibniz' rule, and
absorb coefficients into \(C_n\).  Pair with \(2\nabla^k\Rm\); the Bochner
identity gives the negative gradient term, and variation of the tensor norm
is absorbed into the \(i=0\) summand.
\end{proof}

\begin{lemma}[Main factorial curvature estimate]
\label{lem:chow-ii-main-factorial-curvature-estimate}
\source{chow-et-al-ricci-flow-part-ii:page-0248,chow-et-al-ricci-flow-part-ii:page-0249,chow-et-al-ricci-flow-part-ii:page-0250}
\dcref{ch13:13.29}
\group{Curvature Derivative Estimates}

There is \(C_n<\infty\) such that, if
\(K=\sup_{\mathcal M}|\Rm|(\cdot,0)\), then for every \(m\),
\[
 \sum_{k=0}^m\frac{t^k}{((k+1)!)^2}|\nabla^k\Rm|^2(x,t)
 \leq\frac{K^2}{(1-\tfrac12C_nKt)^2}
\]
on \(0\leq t\leq(4C_nK)^{-1}\).  In particular,
\(|\nabla^k\Rm|(x,t)\leq\frac87(k+1)!Kt^{-k/2}\).
\end{lemma}
\begin{proof}
\sketch
\uses{lem:chow-ii-curvature-derivative-evolution-inequality,thm:noncompact-scalar-pde-ode-comparison}
Set \(A_k=t^{k/2}|\nabla^k\Rm|/(k+1)!\),
\(B_k=t^{(k-1)/2}|\nabla^k\Rm|/k!\),
\(\Phi=\sum_{k=0}^mA_k^2\), and \(\Psi=\sum_{k=1}^mB_k^2\).
Weighted summation, Cauchy--Schwarz, and
\(\sum_{i\geq0}(i+2)^{-2}\leq1\) give
\[
 (\partial_t-\Delta)\Phi
 \leq-(1-2C_nt\Phi^{1/2})\Psi+C_n\Phi^{3/2}.
\]
For \(t\leq(4C_nK)^{-1}\), comparison with
\(y'=C_ny^{3/2}\), \(y(0)=K^2\), and continuity yield the stated bound.
Taking its \(k\)-th summand gives the derivative estimate.
\end{proof}

\begin{example}[Uhlenbeck curvature commutators]
\label{ex:chow-ii-uhlenbeck-curvature-commutators}
\source{chow-et-al-ricci-flow-part-ii:page-0251}
\dcref{ch13:13.30}
\group{Curvature Derivative Estimates}

With Uhlenbeck's time-dependent bundle isometry,
\[
 (\nabla_t-\Delta)\Rm=5\Rm*\Rm,
 \qquad [\nabla_t-\Delta,\nabla_a]\rho
 =R_{abcd}\nabla_b\rho_{cd}.
\]
The factorial estimates have smaller constants than the direct calculation,
but the same factorial growth and analyticity conclusion.
\end{example}
\begin{proof}
\sketch
\uses{lem:chow-ii-main-factorial-curvature-estimate}
Iteration produces only \(\nabla^i\Rm*\nabla^{k-i}\Rm\); moving-metric and
connection terms are absorbed into \(\nabla_t\).  The same Pascal estimates
and weighted \(A_k,B_k\) argument apply with reduced coefficients.  Thus only
\(C_n\), not the factorial order or conclusion, changes.
\end{proof}

\begin{lemma}[Normal-coordinate metric derivatives]
\label{lem:chow-ii-normal-coordinate-metric-derivatives}
\source{chow-et-al-ricci-flow-part-ii:page-0251}
\dcref{ch13:13.31}
\group{Real Analyticity}

Suppose on \(B(p,r)\) that
\(|\nabla^k\Rm|\leq C_1(n)k!r^{-k-2}\).  In normal coordinates at \(p\),
there are \(C_2(n),C_3(n)\) and \(r_1=r_1(r,n)>0\) such that, when
\(d(x,p)\leq\min\{C_2r/\sqrt{C_1},\operatorname{inj}_g(p)\}\),
\[
 \tfrac12(\delta_{ij})\leq(g_{ij})\leq2(\delta_{ij}),\quad
 |\partial^k\Gamma^l_{ij}|\leq C_3k!r_1^{-k-1},\quad
 |\partial^kg_{ij}|\leq C_3k!r_1^{-k}.
\]
\end{lemma}
\begin{proof}
\sketch
Along radial normal geodesics, the Jacobi equation and Grönwall's inequality
give the metric comparison.  Differentiate the Jacobi equation \(k\) times.
The curvature hypothesis and
\(\sum_{i=0}^k\binom{k}{i}i!(k-i)!=(k+1)k!\) give the Christoffel estimate by
induction after shrinking the radius.  Metric compatibility
\(\partial g=\Gamma*g\) gives the metric estimate by a second induction.
\end{proof}

\begin{remark}[Real analyticity in the noncompact case]
\label{rem:chow-ii-noncompact-analyticity}
\source{chow-et-al-ricci-flow-part-ii:page-0252}
\uses{lem:chow-ii-main-factorial-curvature-estimate, lem:chow-ii-normal-coordinate-metric-derivatives}
\dcref{ch13:13.32}
\group{Real Analyticity}


For a complete noncompact bounded-curvature Ricci flow, Shi's local estimates
give \(t^k|\nabla^k\Rm|^2\leq C_k\) on every initial compact time interval.
These bounds justify the global weak maximum principle in the factorial
estimate, so the same normal-coordinate argument proves spatial real
analyticity at every positive time.
\end{remark}

\begin{example}[Problem: attraction of nonnegative Kähler Ricci curvature]
\label{prob:chow-ii-kahler-ricci-attractive-cone}
\dcref{ch13:13.16}
\group{Unpreserved Curvature Conditions}
\source{chow-et-al-ricci-flow-part-ii:page-0234}
Determine whether metrics on
\(\mathbb{CP}^1\hookrightarrow F_k^n\to\mathbb{CP}^{n-1}\), \(0<k<n\),
can collapse the zero section under Kähler--Ricci flow with the mixed-Ricci
\(U(n)\)-invariant soliton on \(L_k^n\) as singularity model.  This would
show that nonnegative Ricci curvature is not attractive.
\end{example}
